\section{Experimental Design}
\label{sec:setup}

\subsection{Data}
\textbf{Diabetes 130-US Hospitals}~\cite{strack2014} is the main tabular benchmark: 30-day readmission from encounter-level features. Cleaning follows Strack \emph{et al.}: encounters discharged to death or hospice are removed and only the first encounter per patient is kept, which leaves \num{diabetesmlpNTrain} training and \num{diabetesmlpNTest} test encounters from \num{diabetesmlpNPatients} patients, \num{diabetesmlpNFeatures} encoded features, and a positive rate of \num{diabetesmlpPrevalence}\%. The public release carries no hospital identifier, so sites are formed from the medical-specialty column by longest-processing-time bucketing with a client-share cap of 15\%: a specialty larger than the cap, including the ``unknown specialty'' group that covers roughly half of the encounters, is split across clients, patient-grouped, before bucketing. The largest client therefore holds \num{diabetesmlpMaxClientShare}\% of the training data and the mean label total-variation distance between clients is \num{diabetesmlpLabelTv}. We call these clients \emph{sites} rather than hospitals. The train/test split and every per-client held-out split are patient-grouped.

\textbf{Real institutions.} Two benchmarks have one client per real institution. \textbf{Fed-Heart-Disease}~\cite{detrano1989,duterrail2022flamby} pools the four UCI Heart Disease collections (Cleveland Clinic, Hungarian Institute of Cardiology, University Hospitals Z\"urich and Basel, VA Long Beach) with the target ``any angiographic disease''; it has \num{realHeartClients} clients and \num{realHeartNtrain} training patients, and missingness indicators are added because two features are mostly missing outside Cleveland. The \textbf{eICU-CRD demo}~\cite{pollard2018,eicudemo} is the open subset of eICU-CRD with its true \texttt{hospitalid}; hospitals with fewer than 80 stays are dropped rather than pooled, so each of its \num{realEicuDemoClients} clients is one hospital, the target is in-hospital mortality, and only admission-time features are used. Both are small; they test whether the conclusions survive real site heterogeneity, not clinical performance. The full credentialed eICU-CRD is supported by the same loader and was not run. \textbf{PneumoniaMNIST}~\cite{yang2023medmnist} is the imaging benchmark: a natively binary $28\times28$ collection, partitioned by Dirichlet label skew ($\alpha=0.3$) with a minimum client size of 50 images. Table~\ref{tab:data} summarises; all dataset statistics in this article are read from the run logs.

\begin{table}[t]
\caption{Dataset characteristics, from the run logs.}
\label{tab:data}
\centering\small
\begin{tabular}{lrr}
\toprule
 & Diabetes 130-US & PneumoniaMNIST\\
\midrule
Modality & tabular EHR & imaging\\
Train / test rows & \num{diabetesmlpNTrain} / \num{diabetesmlpNTest} & \num{medmnistcnnNTrain} / \num{medmnistcnnNTest}\\
Patients (train) & \num{diabetesmlpNPatients} & --\\
Features / channels & \num{diabetesmlpNFeatures} & 1 ($28\times28$)\\
Positive rate (\%) & \num{diabetesmlpPrevalence} & \num{medmnistcnnPrevalence}\\
Sites & 20 (specialty, capped) & 20 (Dirichlet 0.3)\\
Largest site share (\%) & \num{diabetesmlpMaxClientShare} & \num{medmnistcnnMaxClientShare}\\
Label TV (mean) & \num{diabetesmlpLabelTv} & \num{medmnistcnnLabelTv}\\
\bottomrule
\end{tabular}
\end{table}

\subsection{Models and tiers}
Four backbones share one architecture table used by every method: an MLP (hidden 128--128--64), a tabular ResNet (width 128, depth 6), an FT-Transformer (width 64, depth 3, 8 heads)~\cite{gorishniy2021} and a small CNN for imaging. A wide MLP (hidden 1024--1024--512, \num{wideParams} parameters on Diabetes) is used for the link-emulation experiment, so that each upload is several megabytes. Table~\ref{tab:tiers} reports the parameter counts and tier shares of each model as instantiated on the real data.

\begin{table}[t]
\caption{Automatic tier partition on the real data (from the run logs). The fast tier is the private head.}
\label{tab:tiers}
\centering\small
\begin{tabular}{lrrrr}
\toprule
Architecture & Params & Slow (\%) & Medium (\%) & Fast (\%)\\
\midrule
MLP (Diabetes) & \num{diabetesmlpTierTotalParams} & \num{diabetesmlpTierSlowPct} & \num{diabetesmlpTierMediumPct} & \num{diabetesmlpTierFastPct}\\
CNN (PneumoniaMNIST) & \num{medmnistcnnTierTotalParams} & \num{medmnistcnnTierSlowPct} & \num{medmnistcnnTierMediumPct} & \num{medmnistcnnTierFastPct}\\
\bottomrule
\end{tabular}
\end{table}

\subsection{Protocol}
$R=60$ rounds, 20 sites at full participation, one local epoch, batch 64, Adam at $10^{-3}$ with persistent state, inverse-frequency class weighting, $T_w=5$, $K_{\max}=8$, $\rho=0.25$ unless stated. Baselines use their published defaults (FedProx $\mu=0.01$; FedLAMA-style $\tau=5$, $\phi=2$; AdaComm $\tau_0=8$; LAG $\xi=1$); none were tuned on the test data and neither was \tierfed{}. Because personalised methods hold private heads, every method is evaluated on each site's patient-grouped held-out split with that site's own parameters and the per-site metrics are sample-weighted; sites whose held-out split contains a single class cannot yield an AUROC and are excluded identically for every method, with their count reported. A pooled AUROC over all held-out predictions and the global model's AUROC on the untouched test set are reported alongside. Compressors act on each site's update relative to the broadcast with per-client error feedback, and communication is counted from what the compressor emitted for the sites that actually sent. In the fixed multi-rate schedule, as in \tierfed{}, a round in which no tier is due is local-only.

\textbf{Differential privacy.} The matched-privacy comparison uses $C=1$, $\delta=10^{-5}$, targets $\varepsilon^\star\in\{1,2,4,8\}$ at Poisson rate $q=0.5$ and $\varepsilon^\star=4$ at $q=1$, and the four DP-eligible methods FedAvg, FedPer, a fixed period $(K_s,K_m)=(4,4)$ with $T_w=5$, and \tierfed{}-DP with $\lambda=1$ ($\lambda\in\{0.1,10\}$ as sensitivity). Each method's $\sigma$ is calibrated to its own planned release count; final-round AUROC is reported, since best-round selection on held-out data is not covered by the budget.

\textbf{Link emulation.} The analytic cost model (\ref{eq:netcost}) is checked against real transport. Each run's logged per-site upload sizes are replayed, round by round, as concurrent loopback TCP connections, one per transmitting site, whose senders are paced by a token bucket at the profile's uplink rate after a one-way delay of half the RTT; the round ends when the last byte arrives. This is application-layer, trace-driven emulation: it includes TCP segmentation, socket buffering and operating-system scheduling, but not radio scheduling, fading, loss or the downlink, and kernel-level emulation (netem) was not available on the test machine. Time-to-target adds the measured compute time of the slowest site per round.

\subsection{Determinism}
Repeated runs of one configuration are bit-identical on the machine that produced the logs (deterministic kernels, single-threaded reductions, seeded loaders); the reproduction script asserts this before any run. An earlier version of this work lacked this guarantee and identical configurations differed by up to 0.006~AUROC, which is larger than the effect it tested.

\subsection{Statistics}
All comparisons are paired across the ten shared seeds. We report Holm-corrected paired $t$-test $p$-values over the family of all method-versus-FedPer contrasts, BCa bootstrap intervals of the paired difference, and the paired standardised difference $d_z$ with the small-sample correction. Parity is never inferred from a non-significant difference. A one-sided non-inferiority test at margin $\delta_{\mathrm{NI}}=0.01$ is run on \emph{both} final AUROC and final AUPRC,
\begin{equation}
t_{\mathrm{NI}}=\frac{\bar d+\delta_{\mathrm{NI}}}{s_d/\sqrt n},\qquad p=1-\mathcal T_{n-1}(t_{\mathrm{NI}}),
\label{eq:ni}
\end{equation}
and the non-inferiority tests are reported uncorrected and interpreted against the Holm-corrected difference tests. The margin is a reporting convention: 0.01~AUROC is below the between-seed standard deviation of every method in Table~\ref{tab:main}; it has no clinical derivation and we do not claim one. Early-stopped and last-ten-round statistics are reported so that no conclusion depends on the final round alone.
